\documentclass[12pt,a4paper]{article}

% --- Настройка языка и шрифтов для LuaLaTeX ---
\usepackage{fontspec}
\defaultfontfeatures{Ligatures=TeX,Scale=MatchLowercase}

\setmainfont{Alice-Regular}[
    Path=font/,
    Extension=.ttf,
    BoldFont = Alice-Regular,
    ItalicFont = Alice-Regular,
    BoldItalicFont = Alice-Regular,
    BoldFeatures = {FakeBold=1.6},
    ItalicFeatures = {FakeSlant=0.2},
    BoldItalicFeatures = {FakeBold=1.6,FakeSlant=0.2}
]

\newfontfamily\cyrillicfont{Alice-Regular}[
    Path=font/,
    Extension=.ttf,
    BoldFont = Alice-Regular,
    ItalicFont = Alice-Regular,
    BoldItalicFont = Alice-Regular,
    BoldFeatures = {FakeBold=1.6},
    ItalicFeatures = {FakeSlant=0.2},
    BoldItalicFeatures = {FakeBold=1.6,FakeSlant=0.2}
]

\usepackage{polyglossia}
\setmainlanguage{russian}
\setsansfont{TeX Gyre Heros}
\setmonofont{TeX Gyre Cursor}
\newfontfamily\TeXGyreTermes{TeX Gyre Termes}
\newfontfamily\TeXGyrePagella{TeX Gyre Pagella}
\newfontfamily\XCharter{XCharter}

\usepackage{graphicx}        					          % Графика
\usepackage[protrusion=false]{microtype}       	% Микротипографика
\usepackage{amsmath, amsfonts, amssymb, amsthm}	% Математика
\usepackage{braket} 							              % Braket нотация для квантовой механики
\usepackage{unicode-math} 						          % Современный пакет для математики в LuaLaTeX
\usepackage{longtable}       					          % Таблицы с разрывом страниц
\usepackage{tikz}           					          % Рисование
\usepackage[europeanresistors]{circuitikz}		  % Рисование электрических схем в европейском стиле
\usepackage{listings}							              % Для отображения кода
\usepackage{tabularx}       					          % Для таблиц во всю ширину
\usepackage{array}		  						            % Для расширенных возможностей таблиц
\usepackage{float}          					          % Для точного позиционирования фигур
\usepackage{pgfplots}                           % Для построения графиков
\usepackage{caption}                            % Для настройки подписей к рисункам и таблицам
\usepackage{xfp}                                % Для вычислений в LaTeX
\pgfplotsset{compat=1.18}
\usepackage[a4paper,left=20mm,right=20mm,top=20mm,bottom=20mm]{geometry}

\usetikzlibrary{arrows.meta}

% --- Колонтитулы ---
\usepackage{fancyhdr}
\pagestyle{fancy}
\fancyhf{}

\renewcommand{\headrulewidth}{0.4pt}
\renewcommand{\footrulewidth}{0.4pt}
\renewcommand\lstlistingname{Исходный код}

\lstdefinestyle{main}{
	backgroundcolor=\color{gray!8},   
	commentstyle=\color{green!50!black},
	keywordstyle=\color{blue!80!black},
	numberstyle=\small\color{darkgray},
	stringstyle=\rmfamily\color{orange!90!black},
	basicstyle=\ttfamily\small,
	breakatwhitespace=false,         
	breaklines=true,                   
	keepspaces=true,                 
	numbers=left,                    
	numbersep=8pt,                  
	showspaces=false,                
	showstringspaces=false,
	showtabs=false,                  
	tabsize=4
}	

\lstset{style=main}

\newcolumntype{C}{>{\centering\arraybackslash}X}
\renewcommand{\arraystretch}{1.25}

% =========================================================
%         СЛУЖЕБНЫЕ КОМАНДЫ ДЛЯ ВЫЧИСЛЕНИЙ
% =========================================================
\newcommand{\CalcTwo}  [1]{\fpeval{round((#1),2)}}
\newcommand{\CalcThree}[1]{\fpeval{round((#1),3)}}
\newcommand{\CalcFour} [1]{\fpeval{round((#1),4)}}
\newcommand{\CalcFive} [1]{\fpeval{round((#1),5)}}

% ---------------------------------------------------------
% ОБЩИЕ ДАННЫЕ
% ---------------------------------------------------------
\def\VariantN{1}          					 % Номер варианта
\def\LabNumber{2}							  % Номер лабораторной работы
\def\Discipline{Метрология}				     % Название дисциплины
\def\LabTitle{Измерительные сигналы}	 	 % Название лабораторной работы

\title{
	МИНОБРНАУКИ РОССИИ

	федеральное государственное автономное образовательное учреждение

	высшего образования

	«Национальный исследовательский университет
	«Московский институт электронной техники»
	\vspace{2em}
}
\author{
}
\date{}

\begin{document}

\maketitle
\thispagestyle{fancy}

\begin{center}
\Large Отчёт по Лабораторной работе $\LabNumber$\\
по дисциплине «\Discipline»\\
«\LabTitle»
\end{center}

\begin{center}
\Large Вариант \VariantN
\end{center}

\vspace{12em}

\begin{flushright}
	Выполнили студенты группы ЭН-$25$\\
	Андреев И. А.\\
	Анисимов А. С.\\
	Проверил преподаватель\\
	Туринцев К. А.
\end{flushright}

\cfoot{Москва 2026}

\pagebreak

\fancyhead[L]{\LabTitle} 
\fancyhead[R]{\Discipline}
\fancyfoot[C]{\thepage}

\begin{center}
	\Large\textbf{Лабораторная работа №$\LabNumber$} 
\end{center}

\begin{center}
	\LARGE\textbf\LabTitle
\end{center}

\large\textbf{\textit{Цель работы}:} постановка экспериментов по оценке параметров
измерительных сигналов с использованием функционального генератора, цифрового 
мультиметра и цифрового осциллографа с последующей
оценкой погрешностей данных измерений.

\large\textbf{\textit{Необходимое оборудование}:} цифровой мультиметр 
(\textbf{Rigol DM$3058$}), генератор сигналов (\textbf{Rigol DG$1022$Z}), цифровой осциллограф 
(\textbf{Rigol MSO$5102$}), макетная плата, кабели, перемычки, щупы.

\section{Теоретические сведения}



\pagebreak

\begin{center}
	\LARGE{\textbf{Выполнение лабораторной работы и расчетно-графическая часть}}
\end{center}

\fancyhead[L]{Задание $2$} 
\section{Измерение напряжения перменого тока на выходе генератора сигналов}

\fancyhead[L]{Задание $3$}
\section{Измерение методической погрешности из-за подключаемой нагрузки на выход генератора сигналов}\

\fancyhead[L]{Задание $4$}
\section{Измерение частоты среза апериодического звена}

\fancyhead[L]{Задание $5$}
\section{Измерение напряжения переменного тока на выходе ФНЧ}

\fancyhead[L]{Задание $6$}
\section{Измерение напряжения переменного тока на выходе ФВЧ}

\end{document}